\section{Searching over Chains}
\label{sec:search}

\paragraph{Space size.}
With $m = 12$ tactics and chains of length up to three, the ordered search space has $m + m^2 + m^3 = 1{,}884$ chains before intensities are considered. Each chain requires a generation-free pass over \ph{number} documents, a detector evaluation under $L$ keys, and the semantic metrics. Two intensities per step would multiply the space by eight. An exhaustive sweep is possible in principle and too expensive in practice, so we search.

\paragraph{Procedure.}
We use a beam search over depth. At depth one, every tactic is evaluated at its default intensity and scored by the objective in Equation~\ref{eq:objective}. At depth two, each surviving chain is extended with every tactic, including repeats, and re-evaluated. At depth three the same extension applies, truncated to the best $B$ chains. The beam width $B$ is fixed before the search and reported. Every chain is scored on the same document set, so comparisons within a depth are paired. Across depths the comparison is descriptive, since the admissible set changes.

\paragraph{Pruning.}
Two rules keep the budget bounded. A chain whose mean preservation score falls below $\tau - \epsilon$ is dropped, where $\epsilon$ accounts for metric noise. A chain whose removal rate does not exceed that of its best prefix by more than the smallest detectable difference is dropped, because the extra step costs meaning without buying removal. Both rules are logged, so a reader can see which chains were evaluated and which were inferred to be dominated.

\paragraph{Intensity.}
Each tactic has one intensity parameter, normalised to $[0,1]$ where the implementation permits. The main sweep uses \ph{number} settings per tactic. For back-translation the intensity is the number of pivots; for token edits it is the fraction of positions; for paraphrase it is a model-side control such as lexical diversity. We report the chosen intensity with every recommended chain, because a chain that only works at maximum intensity is a different recommendation from one that works at a moderate setting.

\paragraph{Selection and multiple comparisons.}
Searching 1{,}884 chains and reporting the best invites selection bias. We therefore split documents into a search split and a held-out split before any evaluation, choose the chain on the search split, and report its performance on the held-out split. The headline numbers in \S\ref{sec:results} are held-out numbers. Where the budget does not allow a full held-out pass for every chain, we report the top \ph{number} chains on both splits and say so explicitly.

\paragraph{Stopping.}
The search stops at depth three by construction, and earlier when no chain survives pruning. We also report the marginal effect of the third step, which is the quantity that decides whether the three-tactic budget is worth spending at all. If the best pair captures most of the achievable removal, the honest recommendation is a pair.
